\chapter{Monitoreo Activo y Análisis de Sesiones}

Una vez descubiertos los artefactos en el sistema sin que el atacante advirtiera la intervención de los administradores, el equipo pericial enfrentó una disyuntiva táctica fundamental: desconectar inmediatamente el servidor para contener el incidente o mantenerlo en operación bajo vigilancia encubierta para identificar los vectores de comunicación externa, determinar el alcance real de la intrusión y recopilar evidencia material sobre la identidad del operador.

Se tomó la decisión estratégica de establecer un dispositivo de **monitoreo en vivo pasivo**. Este capítulo documenta el seguimiento en tiempo real de las actividades asociadas a la cuenta sospechosa \texttt{duende}, el despliegue de mecanismos de auditoría a nivel de kernel (\textit{Process Accounting}) y el análisis forense de las sesiones interactivas que precedieron a la desconexión física del equipo.

\section{La Decisión Táctica: Vigilancia Pasiva frente a Desconexión Inmediata}

En la metodología forense clásica, la regla general aconseja aislar el sistema comprometido en el menor tiempo posible para evitar la alteración de evidencias volátiles. No obstante, en intrusiones complejas donde coexisten servicios legítimos y artefactos ocultos, una desconexión prematura alerta a los atacantes, quienes pueden proceder al borrado remoto de rastros en otros nodos de la red institucional o migrar sus herramientas a servidores no identificados.

Al optar por la vigilancia pasiva, el equipo de respuesta asumió las siguientes premisas operativas:
\begin{itemize}
    \item Se implementaron réplicas de respaldo de los directorios de usuario para salvaguardar la integridad de los datos académicos.
    \item Se revocaron silenciosamente los privilegios del bit SUID en los ejecutables identificados para bloquear cualquier intento de elevación de privilegios en caso de que el intruso intentara ejecutarlos.
    \item Se instalaron herramientas de auditoría independientes en el sistema operativo para registrar comandos ejecutados al margen de los historiales de \textit{shell} manipulables por el usuario.
\end{itemize}

\section{Perfil Temporal y Patrón de Conexiones (\texttt{wtmp})}

La primera fuente de telemetría analizada fue el archivo de registro de sesiones del sistema (\texttt{/var/log/wtmp}), consultado mediante el comando \texttt{last}. La extracción de los accesos correspondientes al mes de abril de 2004 para el usuario \texttt{duende} permitió reconstruir una cronología rigurosa de actividad.

El Listado~\ref{lst:last_duende_summary} exhibe un extracto representativo de las conexiones registradas, ordenadas de forma cronológica inversa.

\begin{lstlisting}[style=terminal,caption={Extracto representativo de conexiones registradas para la cuenta duende.},label={lst:last_duende_summary}]
duende  pts/0        Sun Apr 27 20:39 - 20:40  (00:00)     victima2.foo.bar.mx
duende  pts/1        Sun Apr 27 15:27 - 15:29  (00:01)     victima2.foo.bar.mx
duende  pts/2        Sat Apr 26 18:20 - 18:39  (00:18)     dsl-198-51-100-164.isp-acceso.net
duende  pts/6        Fri Apr 25 21:09 - 23:21  (02:12)     victima2.foo.bar.mx
duende  pts/7        Fri Apr 25 19:04 - 19:07  (00:03)     victima2.foo.bar.mx
duende  pts/5        Fri Apr 25 14:50 - 14:55  (00:05)     victima2.foo.bar.mx
duende  pts/4        Fri Apr 25 07:10 - 08:10  (00:59)     victima2.foo.bar.mx
duende  pts/0        Wed Apr 23 18:14 - 18:14  (00:00)     dsl-203-0-113-152.isp-acceso.net
duende  pts/3        Tue Apr 22 20:41 - 20:43  (00:02)     victima2.foo.bar.mx
duende  pts/2        Mon Apr 21 16:30 - 17:15  (00:45)     dialup-198-18-10-216.isp-acceso.net
duende  pts/6        Wed Apr 16 12:11 - 12:13  (00:01)     victima2.foo.bar.mx
duende  ftp          Tue Apr 15 03:12 - 03:15  (00:03)     dsl-198-51-100-98.isp-acceso.net
\end{lstlisting}

\subsection{Análisis del Patrón Dual de Accesos}

La correlación pericial de las marcas de tiempo y los orígenes de red reveló un patrón de acceso estrictamente dual:
\begin{enumerate}
    \item \textbf{Accesos locales durante horario institucional:} Durante los días laborales, la cuenta registraba múltiples inicios de sesión breves originados desde el servidor paritario \texttt{victima2.foo.bar.mx}, una máquina ubicada en el mismo edificio que compartía almacenamiento NFS.
    \item \textbf{Accesos externos en horarios no laborales y fines de semana:} Durante las noches, madrugadas y fines de semana (como el sábado 26 y domingo 27 de abril), los accesos provenían invariablemente de conexiones externas conmutadas (dial-up) y enlaces ADSL pertenecientes a proveedores de acceso a Internet residenciales (\texttt{isp-acceso.net}).
\end{enumerate}

Este comportamiento demostró que el operador de la cuenta accedía físicamente a las instalaciones universitarias durante el día e interactuaba de forma remota desde enlaces domésticos fuera del horario habitual de trabajo.

\section{Vigilancia a Nivel de Kernel: Contabilidad de Procesos (\texttt{acct})}

Dado que el intruso había vaciado su historial de comandos en \texttt{.bash\_history}, el equipo de respuesta no podía confiar en los archivos generados en el espacio de usuario. Para capturar los comandos ejecutados sin depender de las herramientas estándar, se habilitó el subsistema de contabilidad de procesos del kernel Linux (\texttt{Process Accounting} o \texttt{acct}), consultado mediante la utilidad \texttt{lastcomm}.

A diferencia de los historiales de \textit{shell}, que son gestionados por el intérprete de comandos y pueden ser modificados o desactivados por el usuario, el subsistema \texttt{acct} registra directamente desde el kernel la terminación de cada proceso (\textit{syscall} \texttt{exit}), almacenando el nombre del ejecutable, el identificador de usuario real (UID), el tiempo de CPU consumido y las banderas de terminación.

El Listado~\ref{lst:lastcomm_duende} presenta un extracto de los comandos capturados por el kernel durante una sesión remota del usuario 500.

\begin{lstlisting}[style=terminal,caption={Registro de ejecución de procesos recuperado mediante lastcomm.},label={lst:lastcomm_duende}]
sshd              SF    500      ??         0.75 secs Sat Apr 26 18:20
tcsh                    500      pts/2      0.10 secs Sat Apr 26 18:20
ls                      500      pts/2      0.02 secs Sat Apr 26 18:21
lsattr                  500      pts/2      0.03 secs Sat Apr 26 18:22
id                      500      pts/2      0.01 secs Sat Apr 26 18:22
su                F     500      pts/2      0.05 secs Sat Apr 26 18:23
ls                      500      pts/2      0.02 secs Sat Apr 26 18:24
cat                     500      pts/2      0.01 secs Sat Apr 26 18:24
gcc                     500      pts/2      1.42 secs Sat Apr 26 18:25
as                      500      pts/2      0.35 secs Sat Apr 26 18:25
collect2                500      pts/2      0.22 secs Sat Apr 26 18:25
ld                      500      pts/2      0.45 secs Sat Apr 26 18:25
\end{lstlisting}

\subsection{Interpretación de las Acciones del Intruso}

La secuencia registrada por el kernel desentrañó con precisión matemática los pasos ejecutados por el atacante:
\begin{itemize}
    \item A las 18:20~h inició sesión vía SSH y ejecutó \texttt{lsattr}, verificando los atributos de inmutabilidad en los archivos que previamente había implantado.
    \item A las 18:23~h ejecutó el comando \texttt{su} para intentar elevar privilegios a \texttt{root}. La bandera \texttt{F} (fallido) confirmó que el intento fue rechazado por el sistema de autenticación.
    \item Entre las 18:24~h y las 18:25~h procedió a compilar nuevo código fuente en lenguaje C invocando directamente el compilador \texttt{gcc}, el ensamblador \texttt{as} y el enlazador dinámico \texttt{ld}.
\end{itemize}

\section{Captura de Variables de Entorno y Sesión de Superusuario}

En una fase posterior de la supervisión, se constató que el usuario logró momentáneamente obtener una sesión de \texttt{root} a través de una brecha residual en un script local. La inspección forense del espacio de memoria del proceso permitió recuperar el volcado de las variables de entorno asociadas a dicha sesión, como se detalla en el Listado~\ref{lst:env_root_session}.

\begin{lstlisting}[style=terminal,caption={Variables de entorno capturadas durante la sesión privilegiada.},label={lst:env_root_session}]
HZ=100
PS1=\h:\w\$ 
USER=root
MAIL=/var/mail/root
LANG=C
LOGNAME=root
SHLVL=1
PWD=/dev/zqft4
PATH=/usr/local/sbin:/usr/local/bin:/sbin:/bin:/usr/sbin:/usr/bin
HOME=/root
_=/usr/bin/env
\end{lstlisting}

La variable \texttt{PWD=/dev/zqft4} aportó una prueba concluyente e irrefutable: la sesión de superusuario no se encontraba operando en el directorio administrativo estándar (\texttt{/root}), sino dentro del directorio clandestino de dispositivos donde residían los bots de IRC. 

\section{Sospecha y Maniobras de Contravigilancia del Atacante}

Hacia finales de abril, los registros revelaron que el intruso comenzó a sospechar que el servidor estaba bajo observación pericial:
\begin{itemize}
    \item Ejecutó comandos de verificación de procesos y consultas repetitivas a los registros del sistema (\texttt{dmesg}, \texttt{tail /var/log/messages}).
    \item Intentó renombrar directorios de trabajo y ejecutar scripts de borrado selectivo para eliminar herramientas intermedias.
    \item Realizó pruebas de conectividad hacia puertos de administración para determinar si se habían instaurado reglas de filtrado en el cortafuegos local.
\end{itemize}

Habiéndose cumplido los objetivos del monitoreo encubierto ---la recopilación de la cadena de mando, los orígenes de red y la confirmación material del vínculo entre la cuenta \texttt{duende} y los ejecutables maliciosos---, la administración procedió formalmente a la desconexión física del servidor y a la clonación en frío de sus discos duros para dar inicio a la fase de análisis forense \textit{post-mortem}.
